\section{Other questions reached by the exported theorems}\label{sec:consequences}

Throughout this section the repository's exported RH and Dirichlet-GRH theorems are used as inputs. This is a mathematical dependency, not an assertion that the present workspace performed their kernel verification. The classical implications listed below are not new theorems of this report.

\subsection{Every abelian Dedekind zeta function over \texorpdfstring{\(\Q\)}{Q}}

\begin{proposition}[Finite products preserve critical-line zero location]\label{prop:product}
Let \(F_1,\ldots,F_r\) be holomorphic in \(0<\Rea s<1\), and suppose each has all its zeros in that strip on \(\Rea s=1/2\). Then the same is true of \(\prod_jF_j\).
\end{proposition}

\begin{proof}
At any point in the strip, a finite product of nonzero complex numbers is nonzero. A zero of the product therefore belongs to at least one factor, whose hypothesis gives the conclusion.
\end{proof}

\begin{corollary}[Abelian Dedekind GRH]\label{cor:abelian}
For every finite abelian extension \(A/\Q\), every zero of \(\zeta_A(s)\) in \(0<\Rea s<1\) lies on \(\Rea s=1/2\).
\end{corollary}

\begin{proof}
Kronecker--Weber places \(A\) in a cyclotomic field. Associate to each character of \(\Gal(A/\Q)\) its primitive Dirichlet character, and let \(X_A\) be this finite set. The exact factorization is
\begin{equation}\label{eq:abelian}
\zeta_A(s)=\prod_{\chi\in X_A}L(s,\chi).
\end{equation}
This equality holds meromorphically; it follows from the Euler factors in the convergence half-plane and continuation. Every factor is holomorphic in the open strip, because its only possible pole is at \(1\). Apply the repository's Dirichlet-GRH theorem and Proposition~\ref{prop:product}.
\end{proof}

The standard factorization and Kronecker--Weber theorem are explicitly stated in Sutherland's course notes \cite{kronecker}. This corollary includes quadratic and cyclotomic fields and their subfields. The primitive-character convention matters for exact equality: writing all factors at one larger imprimitive modulus can remove finite Euler factors. Those factors do not vanish in the open strip, but using the primitive factors avoids the extra correction.

This is a realistic first formalization project. Its new work is the formal factorization interface, not a new estimate for zeros. The included Lean file proves only the generic product and Dirichlet-product parts; it does not pretend that \eqref{eq:abelian} has been formalized there.

\subsection{Real zeros and central signs}

\begin{proposition}[Real nonprincipal characters]\label{prop:positive}
Let \(\chi\) be a real-valued, nonprincipal Dirichlet character. Then
\[
L(\sigma,\chi)>0\quad(\sigma>1/2),
\qquad
L(1/2,\chi)\ge0.
\]
\end{proposition}

\begin{proof}
For real \(\sigma>1\), the absolutely convergent Euler product has positive factors, so \(L(\sigma,\chi)>0\). Because \(\chi\) is real, analytic continuation is real-valued on the real axis. Because it is nonprincipal, it is entire. The exported half-plane theorem says it is nonzero on the connected real interval \((1/2,\infty)\). A continuous real-valued function cannot change sign there without vanishing. It is therefore positive throughout that interval, and continuity at \(1/2\) gives the nonnegative endpoint value.
\end{proof}

This removes real zeros in \((1/2,1)\), including the possible exceptional zeros near \(1\). It does not prove strict positivity at \(1/2\). The nonprincipal assumption is essential: a principal function has a pole at \(1\), so the connectedness argument across that point is unavailable. Likewise, zero location does not assert simplicity or any bound on multiplicities.

\subsection{Quantitative and elementary consequences}

The following are representative established implications with hypotheses matching the exported result. Logarithms are natural. Constants in \(O_\varepsilon\) may depend on \(\varepsilon\).

\begin{longtable}{P{0.24\linewidth}P{0.48\linewidth}P{0.20\linewidth}}
\toprule
Question & Consequence after importing the result & Source / dependency\\
\midrule
\endhead
Least quadratic nonresidue & For prime \(p\ge5\), the least positive nonresidue is \(<(\log p)^2\). & Dirichlet GRH; \cite{lls}, Cor.~1.1.\\[0.3em]
Least prime in a progression & If \(q>3\), \((a,q)=1\), then \(P(a,q)\le(\varphi(q)\log q)^2\). & Dirichlet GRH; \cite{lls}, Cor.~1.2.\\[0.3em]
Lindel\"of bounds & For primitive \(\chi\) of conductor \(q\), \(\lvert L(1/2+it,\chi)\rvert\ll_\varepsilon(q(2+|t|))^\varepsilon\). & Dirichlet GRH and standard analytic estimates; \cite{chandee}.\\[0.3em]
M\"obius cancellation & \(\sum_{n\le x}\mu(n)=O_\varepsilon(x^{1/2+\varepsilon})\) for every \(\varepsilon>0\). & Ordinary RH; \cite{sound}.\\[0.3em]
Primes in short intervals & For every \(x\ge2\), a prime lies in \((x-\frac4\pi\sqrt{x}\log x,x]\). & Ordinary RH; \cite{dudek}, Thm.~2.\\[0.3em]
Divisor-sum inequality & For \(n\ge1\), \(\sum_{d\mid n}d\le H_n+e^{H_n}\log H_n\), with equality only at \(n=1\). & Equivalent to RH; \cite{lagarias}, Thm.~1.1.\\
\bottomrule
\end{longtable}

The least-nonresidue and least-prime statements use the corrected arXiv version of Lamzouri--Li--Soundararajan \cite{lls}. The explicit constants here are useful examples, not assertions of the current optimal bounds. The divisor inequality is an elementary reformulation of the same RH difficulty, rather than an independent analytic breakthrough.

For M\"obius cancellation, the exponent must retain \(\varepsilon>0\). One must not replace it by the much stronger unit-constant bound \(|\sum_{n\le x}\mu(n)|<\sqrt{x}\). The transfer concerns the stated asymptotic estimate. Similarly, Lindel\"of for this family does not imply a corresponding theorem for every higher-degree automorphic \(L\)-function.

These consequences can be used rigorously in ordinary mathematical proofs once their cited implication theorems are supplied. A Lean development needs those bridges as actual formal theorems; calling the result in \code{Main} does not automatically add all of classical analytic number theory to the library.

\subsection{Why several prominent conjectures are not reached}

\paragraph{Artin's primitive-root conjecture.}
The usual conditional proof needs effective splitting estimates in Kummer fields
\[
\Q(\zeta_k,a^{1/k}),\qquad k\text{ squarefree}.
\]
These are generally nonabelian over \(\Q\). For example, the splitting field of \(X^3-2\) is \(\Q(\zeta_3,\sqrt[3]2)\). The cubic is irreducible by Eisenstein's criterion, and its discriminant \(-108\) is not a square, so the Galois group is \(S_3\). Its zeta factorization involves a two-dimensional Artin representation in addition to one-dimensional factors.

Dirichlet GRH controls the latter; it does not supply the missing degree-two zero estimate. Lenstra's formulation makes the needed Dedekind-zeta hypotheses explicit \cite{lenstra}. Thus the familiar assertion that Artin follows from GRH cannot be used here without specifying the larger family of \(L\)-functions.

\paragraph{Legendre's conjecture.}
The displayed short-interval theorem does not place a prime between every pair of consecutive squares. At \(x=(n+1)^2\), its guaranteed interval has length
\[
\frac4\pi(n+1)\log((n+1)^2).
\]
Dividing by the available gap \(2n+1\) gives a ratio asymptotic to \((4/\pi)\log n\), which diverges. The interval extends too far below \(n^2\). This is a precise failure of that deduction, not a proof that RH logically cannot imply Legendre by some other argument.

\paragraph{Twin primes, binary Goldbach, and correlations.}
The repository locates zeros and controls certain one-variable prime sums. It supplies no argument for the shifted prime correlations required for twin primes or the additive correlations required for binary Goldbach. An analogous warning applies to inferring all M\"obius correlations from a bound for \(\sum_{n\le x}\mu(n)\). No solution of those problems is established by this analysis.

\paragraph{Central nonvanishing, simplicity, and BSD.}
An \(L\)-function may have a zero exactly on the critical line with any multiplicity without contradicting GRH. In particular, statements about values or derivatives at the central point require additional information. The occurrence of Galois representations and Selmer groups as tools inside the repository is not a proof of the Birch--Swinnerton-Dyer conjecture or a general special-value formula.
